\documentclass[10pt]{scrartcl}
\usepackage[mdthm,colorsec,noauthor,noasy]{evan}
\usepackage{handout}
\DeclareMathOperator{\im}{range}
\title{MATH 110 --- Week 3}
\subtitle{Homework 3 --- Solutions}
\author{}\date{}\ihead{}
\begin{document}\maketitle
\begin{prob}\textbf{HW3 P1(a).}\par\idea{The preimage inherits the subspace test because $T$ preserves combinations.}\begin{steps}\ii $T0=0\in U$, so zero lies in $T^{-1}(U)$.\ii If $Tv,Tw\in U$, then $T(av+w)=aTv+Tw\in U$. Thus $T^{-1}(U)$ is a subspace.\end{steps}\end{prob}
\begin{prob}\textbf{HW3 P1(b).}\par\idea{Restrict the domain and codomain so the desired intersection is precisely the range.}\begin{steps}\ii Define $R:T^{-1}(U)\to U\cap\im T$ by $Rv=Tv$. The output is in $U$ by the domain condition and in $\im T$ by its form. The map is linear by restriction.\ii Every $w\in U\cap\im T$ is $Tv$ for some $v$, and $Tv\in U$ implies $v\in T^{-1}(U)$; thus $R$ is onto.\ii Since $0\in U$, every element of $\ker T$ belongs to the restricted domain, giving $\ker R=\ker T$. Apply rank--nullity to $R$ on the finite-dimensional subspace $T^{-1}(U)$.\end{steps}\end{prob}
\begin{prob}\textbf{HW3 P1(c).}\par\idea{Use $U=\operatorname{span}(1,x^2)$ to express the derivative condition as a preimage.}\begin{steps}\ii Example 25 gives $\ker D=\operatorname{span}(1)$ and $\im D=\mathcal P_3(\mathbb R)$. Thus $U\cap\im D=U$ has dimension $2$.\ii Part (b) gives $\dim D^{-1}(U)=1+2=3$. Direct integration gives $p=c+bx+(a/3)x^3$, so a confirming basis is $1,x,x^3$.\end{steps}\end{prob}
\begin{prob}\textbf{HW3 P2.}\par\idea{A common kernel and range must have half the domain dimension.}\begin{steps}\ii On $\mathbb R^4$ let $T(x_1,x_2,x_3,x_4)=(x_3,x_4,0,0)$. This is linear coordinatewise; both its kernel and range are $\operatorname{span}(e_1,e_2)$.\ii If such a map existed on $\mathbb R^5$, rank--nullity would give $5=2\dim\ker T$, impossible because dimension is an integer.\end{steps}\end{prob}
\begin{prob}\textbf{HW3 P3.}\par\idea{The constraints define the two outputs, and simple polynomials give both coordinate directions.}\begin{steps}\ii Let $S(p)=(p(1),p'(1))$. Evaluation and differentiation are linear, and its kernel is exactly $U$.\ii $S(1)=(1,0)$ and $S(x-1)=(0,1)$, so $S$ is onto and has rank $2$. Rank--nullity gives $\dim U=5-2=3$.\ii This agrees with HW2: the basis $(x-1)^2,(x-1)^3,(x-1)^4$ has length $3$.\end{steps}\end{prob}
\begin{prob}\textbf{HW3 P4.}\par\idea{Linearity preserves any fixed relation, including its nonzero coefficients.}\begin{steps}\ii If $\sum a_jv_j=0$ with not all $a_j=0$, applying $T$ gives $\sum a_jTv_j=0$ with those same coefficients. Thus the image list is dependent.\ii For the zero map $T:\mathbb R\to\mathbb R$, the independent list $(1)$ maps to $(0)$, which is dependent.\end{steps}\end{prob}
\begin{prob}\textbf{HW3 P5(a).}\par\idea{Each column consists of the monomial coefficients of one basis image.}\begin{steps}\ii $T1=1$, $Tx=x+x^2$, $Tx^2=x^2+2x^3$, $Tx^3=x^3+3x^4$.\ii With the prescribed bases, $A=M(T)=\begin{pmatrix}1&0&0&0\\0&1&0&0\\0&1&1&0\\0&0&2&1\\0&0&0&3\end{pmatrix}$.\end{steps}\end{prob}
\begin{prob}\textbf{HW3 P5(b,c).}\par\idea{Compare output coefficients to check independence of the image columns.}\begin{steps}\ii For $p=a+bx+cx^2+dx^3$, $Tp=a+bx+(b+c)x^2+(2c+d)x^3+3dx^4$. If this is zero, the constant and linear coefficients give $a=b=0$, then the quadratic gives $c=0$, and the cubic gives $d=0$. Thus $\ker T=\{0\}$, nullity $0$.\ii The four images above span the range by linearity, and their independence follows from the same coefficient argument. Hence the range has basis $1,x+x^2,x^2+2x^3,x^3+3x^4$ and dimension $4$.\ii The monomial coordinate vector of $Tp$ is $A(a,b,c,d)^{\mathsf T}$, an arbitrary combination of its four columns. Thus the coordinate range is exactly their span; the polynomials themselves are recovered from their coordinates.\end{steps}\end{prob}
\begin{prob}\textbf{HW3 P6.}\par\idea{In $\mathbb F_2$, each element is its own additive inverse.}\begin{steps}\ii The second equation gives $z=x$. The first becomes $y+w=1$, so $w=1+y$.\ii Choose $x,y\in\{0,1\}$ freely. The four solutions $(x,y,z,w)$ are $(0,0,0,1)$, $(0,1,0,0)$, $(1,0,1,1)$, $(1,1,1,0)$. Substitution verifies both equations modulo $2$.\end{steps}\end{prob}
\begin{prob}\textbf{HW3 P7 (optional).}\par\idea{Translation by a known solution gives a bijection with the kernel.}\begin{steps}\ii Choose $x_0$ with $Ax_0=b$. Then $Ax=b$ iff $A(x-x_0)=0$, so all solutions are uniquely $x_0+u$ with $u\in\ker A$.\ii By rank--nullity, $\dim\ker A=m-\operatorname{rank}A$. In a kernel basis every coefficient is independently $0$ or $1$; uniqueness of basis coordinates gives exactly $2^{m-\operatorname{rank}A}$ solutions.\end{steps}\end{prob}
\begin{prob}\textbf{Exam practice, Question 2(a).}\par\idea{Injectivity converts a relation among images into a relation among the basis vectors.}\begin{steps}\ii If $\sum a_jTv_j=0$, then $T(\sum a_jv_j)=0$. Injectivity gives $\sum a_jv_j=0$ by the zero-kernel criterion.\ii Independence of $v_1,\ldots,v_n$ forces every $a_j=0$.\end{steps}\end{prob}
\begin{prob}\textbf{Exam practice, Question 2(b).}\par\idea{Extend the independent images to a basis so Proposition 21 can define $S$ freely.}\begin{steps}\ii Extend $Tv_1,\ldots,Tv_n$ to a basis $Tv_1,\ldots,Tv_n,w_1,\ldots,w_r$ of $W$ using Week 2.\ii Prescribe $S(Tv_j)=v_j$ and $S(w_k)=0$. Proposition 21 gives a unique linear $S$ on $W$.\ii For $v=\sum a_jv_j$, $STv=\sum a_jS(Tv_j)=\sum a_jv_j=v$. Hence $ST=I_V$.\end{steps}\end{prob}\end{document}
